\documentclass[a4paper]{article}
\usepackage{amsfonts}
\usepackage{amsmath}
\usepackage{amssymb}
\usepackage{graphicx}     

\begin{document}
  
\noindent \large Auteur: Anna Voronovska \\
\noindent \large Studierichting: Informatica\\
\noindent \large Oefeningenreeks 3F nr. 1.2\\

\medskip

\normalsize

$f(x) = \dfrac{2x^2 - x - 4}{x^2 - 2x + 1} = \dfrac{2x^2 - x - 4}{(x - 1)^2}$\\

VA:\\

$(x-1)^2 \neq 0$\\

$x-1 \neq 0$\\

$x \neq 1$\\

Dus VA: $x = 1$\\

Ligging rond de VA:\\

$\lim\limits_{x \to 1^{-}} f(x) = \dfrac{-3}{0+} = -\infty$\\

$\lim\limits_{x \to 1^{+}} f(x) = \dfrac{-3}{0+} = -\infty$\\

Dus de grafiek nadert de VA $x = 1$ langs beide kanten vanonder.\\

HA:\\

$\lim\limits_{x \to \pm\infty} \dfrac{2x^2 - x - 4}{x^2 - 2x + 1} = \lim\limits_{x \to \pm\infty} \dfrac{2x^2}{x^2} = 2$\\

Dus HA: $y = 2$\\

Ligging ten opzichte van de HA:\\

$f(x) - 2 = \dfrac{2x^2 - x - 4}{(x-1)^2} - 2 = \dfrac{2x^2 - x - 4 - 2(x-1)^2}{(x-1)^2} = \dfrac{3x - 6}{(x-1)^2}$\\

$\begin{array}{c|ccccc}
    x         &   & 1     &   & 2 &  \\ \hline
    3x - 6    & - & -     & - & 0 & + \\
    (x - 1)^2 & + & 0     & + & + & + \\ \hline
    f(x) - 2  & - & \vert & - & 0 & +
\end{array}$\\

Als $x \to -\infty$: $f(x) - 2 < 0$, dus grafiek ligt onder de HA\\

Als $x \to +\infty$: $f(x) - 2 > 0$, dus grafiek ligt boven de HA\\


\begin{thebibliography}{999}
%\bibitem{a} Referentie
\end{thebibliography}
\end{document} 
